% Group G12 -- Cloud Computing 100-Day Research Assignment
% Manuscript source. Sections follow Section 14 of the assignment guide.
% If the instructor's Overleaf template uses a different class or section
% order, copy the section bodies into the template and keep references.bib.

\documentclass[11pt,a4paper]{article}
\usepackage[utf8]{inputenc}
\usepackage[T1]{fontenc}
\usepackage[margin=2.5cm]{geometry}
\usepackage{amsmath}
\usepackage{booktabs}
\usepackage{graphicx}
\usepackage{url}
\usepackage[hidelinks]{hyperref}

\title{Thermally Mandated Offloading: Why Quantization Gains Vanish\\ Under Sustained Edge LLM Inference}
\author{
  Faheem \\ 23i-0728 \\ FAST-NUCES Islamabad
  \and
  Irtaza Kazmi \\ 23i-6001 \\ FAST-NUCES Islamabad
}
\date{}

\begin{document}
\maketitle

\begin{abstract}
Edge--cloud placement systems decide where large language model (LLM) inference
runs using a throughput figure measured on a cool, idle device. On consumer GPUs
that figure degrades under sustained load as the device throttles itself, and
temperature, the signal such systems usually monitor, settles long before the
degradation ends. In preliminary single-run measurements on an RTX~3050 laptop
GPU, sustained throughput fell by up to 43\%, and the 21\% cold-start speed
advantage of a 4-bit over an 8-bit quantized model disappeared at thermal
equilibrium. This study characterises sustained performance across model sizes,
quantization levels and two GPUs, tests the mechanism behind the loss of
quantization advantage, builds a predictor of sustained throughput from a short
measurement, and evaluates a placement policy driven by it.
\textit{[Draft abstract for Cutoff~1; rewrite once full results are available.]}

Reproducibility: The code and reproducibility materials for this study are
publicly available at: \url{https://github.com/MFaheemS/Cloud-Project}
\end{abstract}

\section{Introduction}
Running LLMs on local and edge devices is increasingly attractive for cost and
privacy. Deciding \emph{where} a request should run---on the device or in the
cloud---depends on an estimate of how fast the device is. That estimate is
normally taken once, from a short benchmark on cool hardware.

\section{Problem and Motivation}
A consumer GPU running continuously heats up and reduces its own clock speed to
protect itself. It keeps reporting a normal temperature and keeps working, but it
does substantially less work than the benchmark suggested. A placement decision
based on the cold figure can therefore become wrong within minutes, and nothing
in the system detects it.

Quantization, which shrinks a model so it needs less memory and arithmetic, is
one of the main techniques that make edge inference practical. It is almost
always evaluated in short runs, so any interaction between quantization and
sustained thermal behaviour is invisible to current evaluations.

\section{Research Questions and Objectives}
\paragraph{Primary research question.} How does quantization level affect the
throughput a consumer GPU can sustain under continuous LLM inference, and can
steady-state performance be predicted from a short measurement well enough to
make correct edge--cloud placement decisions?

\paragraph{Objectives.}
\begin{enumerate}
  \item Characterise peak and sustained throughput, energy per token, and heating
        and cooling behaviour across model sizes, quantization levels and two GPUs.
  \item Explain the mechanism behind the loss of quantization advantage by
        separating the effect of reduced clock speed from heat itself.
  \item Build and validate a predictor of sustained throughput from a short
        cold-start measurement.
  \item Evaluate a placement policy driven by that predictor against policies
        that assume constant performance or react to temperature.
\end{enumerate}

\section{Literature Review}
% Cutoff 2: expand to ~15-20 papers, organised by theme, with a comparison
% table (problem, method, dataset/setup, metrics, finding, gap).

\paragraph{Edge measurement studies.} Sustained-load benchmarks on mobile, NPU and
laptop GPU platforms confirm heavy throttling under continuous inference
\cite{tummalapalli2026edge}, but test a single model at one quantization level,
are measurement-only, and do not consider placement. Evaluations of quantized
models on edge hardware report energy, accuracy and latency
\cite{husom2025sustainable}, but not how those trade-offs change over a long
session.

\paragraph{Energy benchmarks.} Large cross-GPU energy studies
\cite{fadelargerich2026wattcounts} characterise many models and architectures but
use only 16-bit weights, list quantized models as future work, and measure
steady state without a time axis.

\paragraph{Thermal- and energy-aware management.} Thermal- and power-aware
schedulers exist for datacentres \cite{stojkovic2025tapas}, relying on facility
cooling and power-capping controls that consumer hardware does not expose.
On-device energy managers \cite{zou2026enerinfer} tune a single device rather than
deciding where work should run.

\paragraph{Placement and precision.} Edge--cloud schedulers
\cite{yang2024perllm} assign requests using performance estimates measured in
advance, and dynamic-precision methods \cite{liu2025flexquant} switch quantization
level for accuracy or resource targets without considering heat.

\section{Research Gap}
No existing work measures how quantization level changes a device's ability to
sustain performance under continuous load, or connects that behaviour to the
decision of where inference should run.

\section{Methodology}
\textit{[Written as planned work for Cutoff~1; convert to past tense after the
experiments.]}

\paragraph{Measurement protocol.} Each run cools the GPU to a defined cold start
(below 60\,$^\circ$C), applies 15~minutes of continuous inference, then records
15~minutes of idle
recovery. Temperature, power, clock frequency and throttle status are sampled at
20~Hz via NVML; throughput is recorded per request.

\paragraph{Mechanism.} Key quantization comparisons are repeated with the GPU clock
held at a low fixed frequency on a cool device, to separate reduced compute from
heat.

\paragraph{Predictor and policy.} A curve fitted to the first 60--90~s of each run
predicts steady-state throughput and is validated against the full-length run.
A placement policy using the predictor is compared with a constant-performance
baseline and a temperature-triggered baseline.

\section{Experimental Setup}
\begin{table}[h]
\centering
\caption{Planned experimental matrix.}
\begin{tabular}{ll}
\toprule
Factor & Levels \\
\midrule
Model size   & 0.5B, 1.5B, 3B parameters (Qwen2.5 family) \\
Quantization & Q4\_K\_M, Q8\_0, FP16 \\
Hardware     & RTX 3050 Laptop (4\,GB), RTX 4070 Laptop (8\,GB) \\
Runtime      & Ollama, NVML via \texttt{nvidia-ml-py} \\
Repetitions  & 5 per configuration \\
\bottomrule
\end{tabular}
\end{table}

\section{Results}
\textit{[Preliminary: single runs on the RTX~3050; to be replaced by repeated
measurements.]}

\begin{table}[h]
\centering
\caption{Preliminary peak and sustained throughput (RTX 3050 Laptop, single runs).}
\begin{tabular}{lrrrr}
\toprule
Configuration & Peak (tok/s) & Sustained (tok/s) & Loss & Clock at steady state \\
\midrule
0.5B Q4 & 243.6 & 179.9 & 26.2\% & 1181 MHz \\
1.5B Q4 & 115.8 &  66.2 & 42.9\% &  839 MHz \\
1.5B Q8 &  95.4 &  66.8 & 30.0\% &  832 MHz \\
\bottomrule
\end{tabular}
\end{table}

\section{Discussion}
% To be written after full results.

\section{Limitations}
% Current evidence: single repetitions, one device, uncontrolled ambient
% temperature; FP16 confounded by CPU offload on the 4 GB card.

\section{Future Work}
% To be written.

\section{Conclusion}
% To be written.

\bibliographystyle{plain}
\bibliography{references}

\section*{Authors Profile}
% Short biography of each author, to be written by the authors.

\end{document}
